\documentclass{article}
\usepackage{graphicx} % Required for inserting images

\title{INF266 - Problem Set 3}
\author{Ragnhild Klette, Maya Evensen, Halvard Håvardsrud}
\date{\today}

\begin{document}

\maketitle

\section{Monte Carlo Methods}
\subsection*{Exercise 1}
\begin{enumerate}
    \item Square area: $A_s = 1$. 
    Since the circle has a diameter of 1, the radius is: $r=\frac12$.
    The circle area is defined as: \[A_c= \pi r^2 =\pi(\frac{1}{2})^2 = \frac{\pi}{4}\].
    So the ratio between circle area and square area is: \[X=\frac{A_c}{A_s}=\frac{\frac{\pi}{4}}{1}=\frac{\pi}{4} \]

    \item Monte Carlo estimate $\hat X =$ (fraction of sampled points that fall inside the circle).
    So set $X \approx \hat{X}$: \[\frac{\pi}{4} \approx \hat{X} \rightarrow \pi \approx 4\hat{X}\]
    Therefore the estimator is: \[\hat{\pi}=4\hat{X}\]

    \item  Generated 10,000 i.i.d. samples uniformly distributed in $[0,1]^2$. Done in code.
    
    \item A point is counted as being inside the inscribed circle if $(x_1 - 0.5)^2 + (x_2 - 0.5)^2 \leq 0.25$. From $N= 10,000$ samples, $N_{in} = 7,767$ points fell inside the circle, giving: \[\hat{X} = \frac{7767}{10000} = 0.7767\]
    
    \item Since $\pi = 4\hat{X} = 4$, we estimate: \[\hat{\pi} = 4\hat{X} = 4 \times 0.7767 = 3.1068\]
    The estimate is close to the true value of $\pi \approx 3.1416$, with some deviation due to Monte Carlo sampling error.

    \item Done in code.

    \item Applying the same inside-circle test, we obtained $N_{in,new} = 7021$, therefore: 
    \[\hat{X}_{new} = 0.7021 , \hat{\pi}_{new} = 4\hat{X}_{new} = 2.8084\]
    
    \item (*) In the uniform case, the samples are distributed uniformly over the square, so the fraction of points inside the circle estimates the area ratio $\frac{\pi}{4}$ and $\hat{\pi}$ approximates $\pi$ up to sampling error.

    In the correlated case, the samples are concentrated around the diagonal, meaning the distribution is no longer uniform. Therefore, the empirical fraction inside the circle no longer corresponds to the geometric area ratio. The estimator becomes biased, which explains why $\hat{\pi} = 2.8084$ is significantly lower than $\pi$.

    What we are actually estimating here is the fraction of diamater of the unit circle, that falls inside the diagonal of a unit square. Wich is $\frac{2r}{\sqrt{2}} \approx 0.70710$, close to our estimate of $0.7021$. Ragnhild er kul og Maya er rå. It is worth to mention that some points from our distribution might even fall outside the diagonal of the square. 
    
\end{enumerate}

\section{MC and MCB's}
\subsection*{Exercise 2}

\begin{enumerate}
    \item In MAB we where solving a control problem. This is because we have to learn a policy that maximizes expected cummulative reward by choosing actions. 
    
    \item In MAB we are estimating the expected reward of each arm. This corresponds an action-value function, since we are estimating the value of taking an action rather than the state. In MAB there is no transition between states, thus the action value is given by $q_\pi = E[R | A = a]$.

    \item Because we have statioanaty states in MAB, a trajectory corresponds to a sequence of arms and rewards over time. ie. $\{(A_1, R_1,), ..., (A_n, R_n)\}$. A collection of trajectories correspond to multiple independent runs of the bandit process. Each of wich will generate a different series of actions and rewards. 

    \item The MAB algorithms like $\epsilon$-greedy correspond to MC control. This is because they estimate action values and use them to improve the policy. The MC prediction evaluates a fixed policy, and does not change it the way we do in MAB. 
    
    \item In MAB, a MC prediction would correspond to exploiting a given or earlier learned algorithm. In this situation we effectivly just estimate the expected reward of each arm under the given policy, without doing any policy improvement. 

    \item Both in MAB's and in MC control, $\epsilon$ is used to ensure suffisient exploration. We set this value to ensure that we dont exploit the best policy so far too much. It gives us a better guarantee that all actions are sufficiently explored. In MC control this also means that we need to visit more state-action pairs. In both cases it controlls the relationship between exploring and exploiting. 
    
    \item If we reduxe a standard RL problem solved by MC to a MAB problem we will loose the ability to evaluete a state-action value $q_\pi = E = [R | S = s,  A = a]$. We are essentially loosing the ability to incorporate state transitions and long-term rewards. If we only use MAB, we are trating every RL problem as a stationary single-state problem with immidiate rewards. 
\end{enumerate}
    

\section{Importace sampling}

\begin{enumerate}
    \item $E[R] =  0.3 * 5 + 0.25 * 10 + 0.1 * 20 = 1.5 + 2.5 + 2 = 6$
    \item Done in code 
    \item Done in code. We got 0.57475, close to our expected value of 0.575
    \item The ratio $P(x) = \frac{P_R(x)}{P_S(x)}$ yields: $0\$: 1.166, 5\$ : 0.857, 10\$ : 0.833, 20\$ : 2.0 $. This means that we are twise as likely to get an outcome of 20 \$ from R and less likely to get outcome 5\$ and 10\$.
    \item Done in code. We got 6.0167, close to the true expected value of 6.
    \item Done in code. $\hat{E}[S'] = 0.7595$ in a iteration of our program. The new p gives us $p(0) = 1.4 , p(5) =  1.2, p(10) = 1, p(20) = 0.4$. This means that R and S' is equally likely to get a reward of 10, and that R is less likely to get a reward of 20 than S'. The new $\hat{E} = 6.0248$,  close to our previous value
    \item Done in code  $\hat{E}[S''] = 0.7529$ in a iteration of our program. The new p gives us $p(0) = 350 , p(5) =  0.6, p(10) = 0.5, p(20) = 100$. This means that R is a lot more likely to pick values like 0 and 20, and less likely to pick senter values than s''. The new $\hat{E}$ = 6.0097, close to the original expexted value. 
    \item The results differ heavily in the tendency to pick the outer values 0 and 20. If we get a sample of 0 or 20 this could hugely shift our estimate due to the heavy weights in p wrt. S''. We can say that S' has good coverage but low variance, and that S'' has higher coverage, but a huge variance. This can explain the differences i saw when i at a later run i got $\hat{E} = 5.41$, showing just how much the value can change. 
\end{enumerate}


\section{Q-learning}
\subsection{Exercise 4}
\begin{enumerate}
    \item Q-learning does not rely on importance sampling because its update rule does not estimate expectations under the policy behavior. Instead it uses the max operrator over the next state actions. This makes the update target independent of the the action that was taken. Since there is no probablility-correction needed between the behaviour and target policy, importance ratios are unnessisary. 
\end{enumerate}
\end{document}
